problem:  
Let \((M^n,g,f)\) be a complete noncompact **steady** gradient Ricci soliton satisfying  
\[
\operatorname{Ric} + \nabla^2 f = 0,
\]
and assume that \(M\) has **positive sectional curvature** and that every level hypersurface of the potential \(f\) is compact.  
It is known that in dimension \(n=3\), such a manifold must be rotationally symmetric (the Bryant soliton).  

**Question / Conjecture:**  
For \(n \ge 4\), must every complete noncompact steady gradient Ricci soliton with positive sectional curvature and compact potential level sets be rotationally symmetric, hence isometric (up to scaling) to the higher-dimensional Bryant soliton?  

Alternatively, if this rigidity fails in some higher dimension, can one explicitly construct a genuinely non-rotational steady soliton satisfying these assumptions?  

Why is it a "good" problem:  
This problem formulates a **precise and geometrically natural extension** of a known rigidity theorem from dimension three (Brendle, 2013) to higher dimensions. It asks whether there exist any new geometric phenomena among higher-dimensional steady Ricci solitons beyond the Bryant family.  

- **Profound insight:** It isolates the essential curvature and compactness assumptions under which the soliton potential acts as a convex exhaustion function and probes whether this rigidity mechanism is universal across dimensions. Resolving it would reveal whether the geometric singularity models for the Ricci flow are fundamentally one-dimensional (rotationally symmetric) or if high-dimensional phenomena occur.  

- **Bridge across fields:** The question stands at the intersection of **geometric analysis** (Ricci flow, soliton equations), **Riemannian geometry** (curvature positivity and convex functions), and **topology** (structure and asymptotic shape of noncompact manifolds with convex exhaustion). Any advance may necessitate new analytic tools (e.g., curvature pinching estimates, stability analysis) and could influence our understanding of **singularity formation under Ricci flow**.  

- **Extreme simplicity:** The statement can be explained in one sentence—“Are all steadily shrinking curvature solitons with positive curvature and compact level sets necessarily rotationally symmetric?”—yet it lies at the frontier of geometric analysis. It transforms a central feature of Ricci flow singularities into a direct and accessible geometric question, exemplifying the “narrow entrance, great depth” ideal of a truly good mathematical problem.